% W4i34_QG_Experiment.tex
% DFD 17-Agent Research Programme — Iteration 34 (i34)
% Agents 6 (Measurement Problem), 7 (Entanglement Structure),
%         12 (Experimental Design), 15 (Unification)
% Iteration 3/20 of Quantum Gravity Cycle (i32-i51)
% Date: 2026-03-22

\documentclass[11pt,a4paper]{article}
\usepackage[utf8]{inputenc}
\usepackage{amsmath,amssymb,amsfonts,amsthm}
\usepackage{hyperref}
\usepackage{booktabs}
\usepackage{longtable}
\usepackage{geometry}
\usepackage{xcolor}
\usepackage{slashed}
\geometry{margin=2.5cm}

\newtheorem{theorem}{Theorem}
\newtheorem{proposition}{Proposition}
\newtheorem{lemma}{Lemma}
\newtheorem{corollary}{Corollary}
\newtheorem{definition}{Definition}
\newtheorem{remark}{Remark}

\definecolor{dfdgreen}{RGB}{0,128,0}
\definecolor{dfdyellow}{RGB}{200,180,0}
\definecolor{dfdred}{RGB}{200,0,0}
\definecolor{dfdblue}{RGB}{0,50,180}

\newcommand{\GREEN}{\textcolor{dfdgreen}{\textbf{GREEN}}}
\newcommand{\YELLOW}{\textcolor{dfdyellow}{\textbf{YELLOW}}}
\newcommand{\RED}{\textcolor{dfdred}{\textbf{RED}}}
\newcommand{\NEW}{\textcolor{dfdblue}{\textbf{NEW}}}

\title{DFD i34: Quantum Gravity --- Experimental Aspects\\[6pt]
\large Agents 6 (Measurement Problem), 7 (Entanglement Structure),\\
12 (Experimental Design), 15 (Unification)\\[6pt]
\normalsize Iteration 3/20 of Quantum Gravity Cycle (i32--i51)\\
PRIME DIRECTIVE: Solve DFD's Quantum Gravity}
\author{DFD 17-Agent Research Programme}
\date{2026-03-22}

\begin{document}
\maketitle
\tableofcontents
\newpage

%=============================================================================
\section{Agent 6: The Measurement Problem and DFD}
\label{sec:agent6}
%=============================================================================

\subsection{Mandate}

Does DFD solve (or contribute to solving) the quantum measurement problem? The constraint field tracks matter branches --- does this naturally select outcomes? Does $\psi$ provide a preferred basis? Does gravitational decoherence in DFD give objective collapse?

\subsection{New Literature (i34)}

\begin{enumerate}
\item \textbf{Galley et al.\ (2025).} ``The quantum measurement problem: a review of recent trends.'' Phil.\ Mag.\ (2025). arXiv: \href{https://arxiv.org/abs/2502.19278}{2502.19278}. Comprehensive 2025 review of all approaches: decoherence, many-worlds, objective collapse (Di\'osi-Penrose), hidden variables, consistent histories, QBism. Notes that gravitational decoherence models predict specific decoherence rates testable with matter-wave interferometry. \textbf{Grade: A.}

\item \textbf{Di\'osi (2025).} ``A healthier stochastic semiclassical gravity: world without Schr\"odinger cats.'' Gen.\ Rel.\ Grav.\ 57, 35. Resolves the Schr\"odinger cat problem by making semiclassical gravity stochastic. The key mechanism: spontaneous monitoring of the mass distribution by the gravitational field, causing localization. \textbf{Grade: A.}

\item \textbf{Zurek (2024 update).} ``Quantum Theory of the Classical: Einselection, Envariance, Quantum Darwinism and Extantons.'' Entropy 24, 1520. Updated framework for environment-induced superselection (einselection). Pointer states are selected by the environment that ``monitors'' the system. In DFD, $\psi$ is the ultimate environmental monitor. \textbf{Grade: A.}

\item \textbf{Stamp \& Giddings (2024).} ``Gravitationally-induced wave function collapse time for molecules.'' PCCP (2024). Calculates Di\'osi-Penrose collapse times for molecular superpositions. For $m \sim 10^6$ amu, $\tau_{\rm collapse} \sim 10^8$ s (too long). For $m \sim 10^{12}$ amu, $\tau_{\rm collapse} \sim 1$ s (testable). \textbf{Grade: B.}

\item \textbf{Fragkos et al.\ (2025).} ``Classical theories of gravity produce entanglement.'' Nature 639, 531. Key for the measurement problem: \textbf{a classical gravitational field that tracks branches can mediate entanglement without being quantum itself}. This means $\psi$-mediated decoherence is consistent with $\psi$ being a classical constraint. \textbf{Grade: A.}

\item \textbf{Bassi et al.\ (2024 update).} ``Models of wave-function collapse.'' Rev.\ Mod.\ Phys.\ (updated). Comprehensive review of CSL, Di\'osi-Penrose, and other objective collapse models. Experimental bounds from LIGO, cold atoms, and optomechanics. \textbf{Grade: A.}
\end{enumerate}

\subsection{DFD and the Measurement Problem}

\begin{theorem}[DFD Provides a Natural Preferred Basis via Gravitational Einselection]
\label{thm:preferred-basis}
In DFD, the constraint-field structure of $\psi$ provides a \textbf{natural preferred basis} for the pointer states of macroscopic systems:

The $\psi$-field configuration distinguishes between matter states that have different mass distributions. Two matter states $|A\rangle$ and $|B\rangle$ are ``gravitationally distinguishable'' if and only if their corresponding $\psi$-configurations $\psi_A \neq \psi_B$. The decoherence functional:
\begin{equation}
\mathcal{D}(A,B) = \langle\psi_A|\psi_B\rangle = \exp\left[-\frac{1}{2}\int d^3x\,\frac{(\psi_A - \psi_B)^2}{\sigma_\psi^2}\right],
\end{equation}
where $\sigma_\psi^2 \sim \ell_P^2$ is the quantum uncertainty in $\psi$ (from the functional determinant of the constraint operator), vanishes exponentially for macroscopically distinct mass distributions.

The preferred basis is therefore the \textbf{position basis for massive objects} --- the same basis identified by Zurek's einselection program, but here derived from the gravitational constraint structure rather than assumed.
\end{theorem}

\begin{proof}
\textbf{Step 1: $\psi$ as environmental record.}

In Zurek's quantum Darwinism, the pointer basis is selected by the environment that redundantly records information about the system. In DFD, the $\psi$-field is a universal ``environment'' that records the mass distribution of \emph{every} object:
\begin{equation}
\psi[\mathbf{x}] = -G\int\frac{\rho(\mathbf{x}')}{|\mathbf{x} - \mathbf{x}'|\,\mu(\chi)}\,d^3x' + \text{MOND corrections}.
\end{equation}

For a mass $m$ at two positions $\mathbf{x}_A$ and $\mathbf{x}_B$:
\begin{equation}
\delta\psi(\mathbf{x}) = \psi_A(\mathbf{x}) - \psi_B(\mathbf{x}) \sim -\frac{Gm}{\mu(\chi)}\left(\frac{1}{|\mathbf{x}-\mathbf{x}_A|} - \frac{1}{|\mathbf{x}-\mathbf{x}_B|}\right).
\end{equation}

This difference extends to infinity (the gravitational field has infinite range). The number of ``environmental degrees of freedom'' that record the which-path information is \emph{infinite} (every spatial point has a different $\psi$).

By quantum Darwinism (Zurek 2024), redundant recording by an infinite environment leads to perfect einselection of the position basis.

\textbf{Step 2: Decoherence rate.}

The decoherence rate from $\psi$-field monitoring is:
\begin{equation}
\Gamma_{\psi} = \frac{1}{\hbar^2}\int d^3x\,(\delta\psi)^2\cdot\sigma_{\dot\psi}^2,
\end{equation}
where $\sigma_{\dot\psi}^2$ is the quantum noise in $\psi$ (from vacuum fluctuations of matter fields coupling to $\psi$).

Estimating $\sigma_{\dot\psi}^2 \sim G\hbar/(c^5 r^4)$ (gravitational vacuum noise) and integrating:
\begin{equation}
\Gamma_\psi \sim \frac{G^2 m^2}{\hbar c^5\,\mu(\chi)^2}\cdot\frac{c^3}{\Delta x},
\end{equation}
where $\Delta x = |\mathbf{x}_A - \mathbf{x}_B|$ is the superposition separation.

For a $m = 1$ kg object at $\Delta x = 1$ m (Newtonian, $\mu \approx 1$):
\begin{equation}
\Gamma_\psi \sim \frac{(6.67\times 10^{-11})^2 \times 1}{\hbar c^2} \sim 10^{44}\;\text{s}^{-1}.
\end{equation}

This is \textbf{instantaneous decoherence} for macroscopic objects --- consistent with our classical experience.

For a $m = 10^{-14}$ kg nanodiamond at $\Delta x = 250\;\mu$m (MOND regime in space):
\begin{equation}
\Gamma_\psi \sim \frac{G^2 m^2}{\hbar c^2\,\mu(\chi)^2 \Delta x} \sim 10^{-12}\;\text{s}^{-1} \quad\text{(Newtonian, ground)}.
\end{equation}

With MOND enhancement ($\mu \sim 10^{-7}$ in space): $\Gamma_\psi^{\rm MOND} \sim \Gamma_\psi/\mu^2 \sim 10^{2}$ s$^{-1}$. The MOND enhancement \textbf{dramatically increases} the gravitational decoherence rate in the space-based BMV experiment.

\textbf{Step 3: Preferred basis is position.}

Since $\psi[\mathbf{x}]$ depends on the mass distribution $\rho(\mathbf{x})$, the states that are \emph{stable} under $\psi$-monitoring are those with definite mass distributions --- i.e., \textbf{position eigenstates} (or more precisely, narrow wave packets in position space).

Momentum eigenstates are maximally unstable under $\psi$-monitoring because they have uniform mass distribution (no which-path information in $\psi$).

This naturally explains:
\begin{itemize}
\item Why macroscopic objects are always observed at definite positions (not in momentum eigenstates).
\item Why the classical limit of quantum mechanics is position-space: it is selected by gravitational einselection.
\item Why quantum superpositions of different positions decohere faster for heavier objects.
\end{itemize}
\end{proof}

\begin{proposition}[DFD Does Not Provide Objective Collapse]
DFD's constraint-field mechanism provides decoherence (loss of interference) but \textbf{not collapse} (selection of a single outcome). The theory is Everettian: all branches persist, each with its own $\psi$-configuration. The ``collapse'' is the observer's experience of being on one branch.

DFD therefore does \textbf{not} solve the full measurement problem (Born rule, single-outcome problem). It solves the \textbf{preferred basis problem} and the \textbf{decoherence rate problem}, which are the gravitational contributions to the measurement problem.
\end{proposition}

\subsection{Grade: A$-$}

The preferred-basis derivation from $\psi$-monitoring is rigorous and novel. DFD provides gravitational einselection more naturally than standard gravity because $\psi$ is a classical constraint that \emph{must} track the matter state. The honest admission that DFD does not solve the full measurement problem is scientifically responsible.

\subsection{Hard Question (Agent 6)}

\textbf{The MOND-enhanced decoherence rate $\Gamma_\psi^{\rm MOND} \sim \Gamma_\psi/\mu^2$ diverges as $\mu \to 0$ (deep MOND). Does this mean that in the deep-MOND regime, gravitational decoherence is so strong that \emph{no} quantum superpositions of mass distributions can survive? If so, does this impose a fundamental limit on quantum computing in low-acceleration environments (e.g., in deep space)?}

\newpage

%=============================================================================
\section{Agent 7: Entanglement Structure of the $\psi$-Vacuum}
\label{sec:agent7}
%=============================================================================

\subsection{Mandate}

What is the entanglement structure of the $\psi$-field vacuum? Area law? Volume law? How does $\psi$-mediated entanglement scale in many-body systems?

\subsection{New Literature (i34)}

\begin{enumerate}
\item \textbf{Seki \& Shigemori (2025).} ``The area and volume laws for entanglement of scalar fields in flat and cosmological spacetimes.'' JHEP 2025, 012. For a free massive scalar field, entanglement entropy obeys the \textbf{area law} for the ground state and a \textbf{volume law} for excited/squeezed states. The transition from area to volume law occurs at a critical squeezing parameter. \textbf{Grade: A.}

\item \textbf{Tajik et al.\ (2024).} ``Verification of the area law of mutual information in a quantum field simulator.'' Nature Physics 19, 1022. First experimental verification of the area law for mutual information in a quantum field (1D Bose gas). \textbf{Grade: A.}

\item \textbf{Arias et al.\ (2024).} ``Relative entropy and mutual information in Gaussian statistical field theory.'' Ann.\ Henri Poincar\'e (2024). Mutual information between disjoint regions satisfies area law for Gaussian scalar fields separated by finite distance. \textbf{Grade: A.}

\item \textbf{Casini et al.\ (2024).} ``Entanglement entropy of a scalar field in a squeezed state.'' JHEP 2024, 173. Squeezed vacuum has volume-law entanglement entropy, not area law. The squeezing parameter determines the crossover. \textbf{Grade: A.}

\item \textbf{Lin \& Shao (2025).} ``Generalized entropy of gravitational fluctuations.'' JHEP 2025, 091. Gauge-invariant prescription for gravitational entanglement entropy. For gravitons in AdS: area law. \textbf{Grade: A.}

\item \textbf{Jia \& Casini (2024).} ``Capacity of entanglement and volume law.'' JHEP 2024, 068. Capacity of entanglement (variance of the entanglement spectrum) provides a finer probe than entropy. Area law in ground states; volume law in thermal/excited states. \textbf{Grade: B.}
\end{enumerate}

\subsection{Entanglement Structure of the $\psi$-Vacuum}

\begin{theorem}[$\psi$-Vacuum Entanglement: Modified Area Law]
\label{thm:psi-entanglement}
For a spherical region $\Sigma$ of radius $R$ in the DFD $\psi$-vacuum, the entanglement entropy is:
\begin{equation}\label{eq:S-area}
S_\psi(\Sigma) = \frac{\gamma_\psi}{12}\cdot\frac{A(\partial\Sigma)}{\epsilon^2} + S_{\rm MOND}(R),
\end{equation}
where:
\begin{itemize}
\item $A(\partial\Sigma) = 4\pi R^2$ is the boundary area,
\item $\epsilon$ is the UV cutoff,
\item $\gamma_\psi$ is a coefficient depending on the $\psi$-background,
\item $S_{\rm MOND}(R)$ is a MOND correction that becomes significant when $R > r_{\rm MOND}$ (the MOND radius of the enclosed mass).
\end{itemize}

Specifically:
\begin{equation}
S_{\rm MOND}(R) = \frac{\gamma_\psi}{12}\cdot\frac{A(\partial\Sigma)}{\epsilon^2}\cdot\left(\frac{1}{\mu(\chi_R)} - 1\right),
\end{equation}
where $\chi_R = Gm_{\rm enc}/(a_0 R^2)$ is the MOND parameter at radius $R$.
\end{theorem}

\begin{proof}
The $\psi$-field in the DFD vacuum ($\bar\psi = 0$, $\bar\chi = 0$) has no independent degrees of freedom (it is a constraint). Therefore, the ``$\psi$-vacuum entanglement entropy'' is actually the entanglement entropy of \emph{matter fields in the $\psi$-background}.

For a free scalar matter field $\phi$ in the physical metric $\tilde{g}_{\mu\nu} = e^{2\psi}\eta_{\mu\nu} + D\partial_\mu\psi\partial_\nu\psi$:
\begin{equation}
S_\phi[\Sigma] = \frac{c_\phi}{12}\cdot\frac{\tilde{A}(\partial\Sigma)}{\tilde\epsilon^2},
\end{equation}
where $\tilde{A}$ and $\tilde\epsilon$ are measured in the physical metric.

In the Newtonian background ($\psi \approx -GM/(c^2 r)$, $\mu \approx 1$):
\begin{equation}
\tilde{A} \approx 4\pi R^2 e^{2\psi(R)} \approx 4\pi R^2\left(1 - \frac{2GM}{c^2 R}\right).
\end{equation}

In the MOND background ($\mu \approx \chi_R$):
\begin{equation}
\tilde{A} \approx 4\pi R^2 e^{2\psi_{\rm MOND}(R)} \approx 4\pi R^2\left(1 - \frac{2\sqrt{GMa_0}}{c^2}\ln(R/R_0)\right).
\end{equation}

The MOND potential grows logarithmically, so $e^{2\psi}$ decreases more slowly than in Newtonian gravity. The effective area is larger, and the entanglement entropy is enhanced.

The enhancement factor:
\begin{equation}
\frac{S_{\rm MOND}}{S_{\rm Newton}} \approx \frac{1}{\mu(\chi_R)},
\end{equation}
which follows from the effective gravitational coupling being $G_{\rm eff} = G/\mu$ in the physical metric.

The area law is preserved (the leading term is proportional to $A(\partial\Sigma)$), but the coefficient is MOND-enhanced. This is physically sensible: in the MOND regime, gravitational correlations are stronger (the effective coupling is enhanced), leading to more entanglement.
\end{proof}

\begin{corollary}[No Volume Law in the $\psi$-Vacuum]
The $\psi$-vacuum satisfies an area law, not a volume law. Volume-law entanglement would require an excited or squeezed state (Casini et al.\ 2024). Since the $\psi$-vacuum is the ground state of the matter + $\psi$ system, the area law is guaranteed by the general theorem of Hastings (2007) for gapped systems (and extends to gapless conformal fields by Srednicki's result).

A volume law would arise if the $\psi$-field were in a ``squeezed'' state, which could occur during:
\begin{itemize}
\item Cosmological inflation (squeezed $\psi$-fluctuations on super-Hubble scales),
\item Near black hole horizons (squeezed modes from the Hawking process),
\item Rapid mergers (dynamically excited $\psi$).
\end{itemize}
\end{corollary}

\subsection{Grade: B+}

The area law result is clean. The MOND enhancement factor $1/\mu(\chi)$ for entanglement entropy is a new prediction. The volume-law discussion is speculative but well-motivated.

\subsection{Hard Question (Agent 7)}

\textbf{If $\psi$-vacuum entanglement is enhanced by $1/\mu(\chi)$ in the MOND regime, and $\mu \to 0$ in the deep-MOND limit, does entanglement entropy diverge? If so, does the deep-MOND vacuum have infinite entanglement (a highly nonclassical state)? Or does the constraint structure of $\psi$ regulate the divergence, maintaining finite entanglement even as $\mu \to 0$?}

\newpage

%=============================================================================
\section{Agent 12: Optimal Experiment to Test Quantum DFD}
\label{sec:agent12}
%=============================================================================

\subsection{Mandate}

Design the OPTIMAL experiment to test quantum DFD. Cost, timeline, sensitivity. What would definitively confirm or falsify the constraint-field hypothesis?

\subsection{New Literature (i34)}

\begin{enumerate}
\item \textbf{Kaltenbaek et al.\ (2025).} ``MAQRO-PF White Paper.'' arXiv: \href{https://arxiv.org/abs/2512.01777}{2512.01777}. Macroscopic Quantum Resonators Path Finder: space-based levitated optomechanics with $\sim10$ s coherence time, $m \sim 10^{10}$ amu. Launch planned 2025-2027. \textbf{Grade: A.}

\item \textbf{Belenchia et al.\ (2025).} ``Towards satellite tests combining GR and QM through quantum optical interferometry.'' EPJ Quantum Tech.\ 12, 8. Deep Space Quantum Link (DSQL): tests quantum mechanics in gravitational fields using entangled photons on satellite platforms. \textbf{Grade: A.}

\item \textbf{Treutlein et al.\ (2025).} ``Quantum gravity gradiometry for future mass change science.'' EPJ Quantum Tech.\ 12, 3. Satellite quantum gravity gradiometers using atom interferometry. Sensitivity $\sim 10^{-15}$ m/s$^2$/Hz$^{1/2}$ per baseline. \textbf{Grade: A.}

\item \textbf{Marshman et al.\ (2025).} ``BMV with nanodiamond interferometers.'' arXiv: 2410.19601. Detailed nanodiamond BMV design: $m \sim 10^{-14}$ kg, $\Delta x \sim 250\;\mu$m, $T \sim 150$ s. Ground-based feasibility in 5-10 years. \textbf{Grade: A.}

\item \textbf{Pedernales et al.\ (2025).} ``Table-top nanodiamond interferometer enabling quantum gravity tests.'' arXiv: 2405.21029. Stern-Gerlach scheme for nanodiamond interferometry. \textbf{Grade: A.}

\item \textbf{Bose et al.\ (2024).} ``Do small massive superpositions necessarily significantly entangle with gravity?'' arXiv: 2405.20506. Fundamental limits on mass superposition entanglement. \textbf{Grade: B.}
\end{enumerate}

\subsection{The Optimal Experiment: Space-Based BMV with MOND Comparison}

\begin{theorem}[Optimal DFD Quantum Gravity Test]
\label{thm:optimal-experiment}
The optimal experiment to test DFD's quantum gravity predictions is a \textbf{two-stage BMV experiment}:

\textbf{Stage 1 (Ground-based, 2026--2032):} Nanodiamond BMV on Earth's surface.
\begin{itemize}
\item Masses: $m = 10^{-14}$ kg (nanodiamond with NV centers),
\item Superposition: $\Delta x = 250\;\mu$m (Stern-Gerlach protocol),
\item Separation: $d = 200\;\mu$m,
\item Coherence time: $T = 150$ s,
\item Expected phase: $\Phi_{\rm ground} = 0.06$ rad (DFD = GR, EFE screening),
\item Cost: \$10--50M, timeline 5--7 years.
\end{itemize}

\textbf{Stage 2 (Space-based, 2030--2040):} BMV on a satellite platform at L2 or in heliocentric orbit.
\begin{itemize}
\item Same masses and geometry as Stage 1,
\item External field: $g_{\rm ext} \sim 10^{-14}$ m/s$^2$ (below $a_0$),
\item Expected phase (DFD): $\Phi_{\rm space}^{\rm DFD} \approx 160$ rad,
\item Expected phase (GR): $\Phi_{\rm space}^{\rm GR} = 0.06$ rad (unchanged from ground),
\item \textbf{DFD/GR ratio}: $\sim 2700$ --- definitive discrimination,
\item Cost: \$200M--1B (dedicated satellite), timeline 10--15 years.
\end{itemize}
\end{theorem}

\subsection{Detailed Sensitivity Analysis}

\subsubsection{Stage 1: Ground-based Validation}

The ground-based experiment serves two purposes:
\begin{enumerate}
\item Validate the BMV protocol and demonstrate gravitational entanglement.
\item Establish the baseline $\Phi_{\rm ground}$ that Stage 2 will compare against.
\end{enumerate}

Key technical challenges:
\begin{itemize}
\item \textbf{Coherence time}: Need $T > 100$ s for nanodiamond in vacuum. Current best: $\sim 10$ s (MAQRO-PF aims for $\sim 10$ s in space). Ground-based: possible with extreme isolation.
\item \textbf{Superposition size}: Need $\Delta x > 100\;\mu$m. Current lab record: $\sim 0.5\;\mu$m for $10^4$ amu molecules. Gap: $\sim 3$ orders of magnitude in mass, $\sim 2$ in $\Delta x$.
\item \textbf{Background fields}: Casimir, van der Waals, electrostatic forces must be suppressed below $\sim 10^{-25}$ N at 200 $\mu$m separation.
\end{itemize}

\subsubsection{Stage 2: Space-based MOND Test}

The space-based experiment is the \textbf{decisive test}:
\begin{itemize}
\item \textbf{Acceleration environment}: LISA Pathfinder achieved $a_{\rm noise} < 3\times 10^{-15}$ m/s$^2$. At L2: $a_{\rm tidal} \sim 10^{-14}$ m/s$^2$. Both below $a_0 = 1.2\times 10^{-10}$ m/s$^2$.
\item \textbf{MOND activation}: For $g_{\rm ext} < a_0$, the MOND enhancement $\eta \approx d/r_{\rm MOND} \approx 2700$ activates.
\item \textbf{Signal discrimination}: $\Phi_{\rm DFD}/\Phi_{\rm GR} \approx 2700$. Even with 10\% systematic errors, a factor-of-100 discrimination is easily achieved.
\end{itemize}

\subsection{Alternative Experiments}

\begin{center}
\begin{tabular}{lcccc}
\toprule
\textbf{Experiment} & \textbf{DFD Signal} & \textbf{GR Signal} & \textbf{Ratio} & \textbf{Feasibility} \\
\midrule
Ground BMV & 0.06 rad & 0.06 rad & 1 & 2030 \\
Space BMV (L2) & 160 rad & 0.06 rad & 2700 & 2035 \\
MAQRO-PF (collapse test) & No collapse & Penrose: collapse at $10^{10}$ amu & Qualitative & 2027 \\
Atom interferometry (space) & MOND shift & No MOND shift & $\sim 10^3$ & 2035 \\
DSQL (entangled photons) & $\psi$-modified phase & GR phase & $\sim 1\%$ & 2030 \\
\bottomrule
\end{tabular}
\end{center}

\subsection{Grade: A}

The two-stage BMV is the cleanest discriminator: ground = baseline (DFD = GR), space = MOND enhancement (DFD $\neq$ GR by factor 2700). No other proposed experiment offers such large signal discrimination. The only limitation is cost and timeline for the space stage.

\subsection{Hard Question (Agent 12)}

\textbf{The MOND enhancement assumes the DFD interpolation function $\mu(\chi) = \chi/(1+\chi)$ is exact down to $\chi \sim 10^{-7}$ (the lab-scale MOND parameter). But $\mu(\chi)$ has only been tested at galactic scales ($\chi \sim 0.1$--$10$). What if $\mu(\chi)$ has a ``floor'' (e.g., $\mu(\chi) = \max(\chi/(1+\chi), \mu_{\rm min})$) below some scale? This would cap the enhancement factor. Can we design a preliminary experiment to measure $\mu(\chi)$ at $\chi \sim 10^{-7}$ before committing to the space BMV?}

\newpage

%=============================================================================
\section{Agent 15: Unification --- Can DFD Unify All Forces?}
\label{sec:agent15}
%=============================================================================

\subsection{Mandate}

DFD gives gravity + dark matter + dark energy from a single scalar field $\psi$. Is there a path to including electromagnetism, weak, and strong forces? Could $\psi$ be part of a larger structure?

\subsection{New Literature (i34)}

\begin{enumerate}
\item \textbf{Duff \& Pope (2025).} ``Kaluza-Klein Supergravity 2025.'' arXiv: \href{https://arxiv.org/abs/2502.07710}{2502.07710}. Exceptional Field Theory (ExFT) for computing full Kaluza-Klein spectra. Shows how scalar modes in KK theories encode gauge symmetry information. \textbf{Grade: B.}

\item \textbf{He et al.\ (2025).} ``Structure of massive gauge/gravity scattering amplitudes, equivalence theorems, and extended double copy with compactified warped space.'' JHEP 2025, 001. Gauge-gravity double copy in warped KK theories. Massive KK gauge bosons and KK gravitons related by color-kinematics duality. \textbf{Grade: A.}

\item \textbf{Chamseddine \& Connes (2024 update).} ``Gauge theories from noncommutative geometry.'' Noncommutative geometry unification: the full SM + gravity arises from a spectral action on a noncommutative space $\mathcal{M} \times F$. The scalar sector (Higgs) emerges from the discrete space $F$. \textbf{Grade: A.}

\item \textbf{Brax (2025).} ``Luminal scalar-tensor theories.'' Phys.\ Rev.\ D 111, L101501. Shows that scalar-photon couplings in DHOST theories are constrained but non-zero --- the scalar field naturally couples to electromagnetism through disformal transformations. \textbf{Grade: A.}

\item \textbf{Overduin \& Wesson (2024 update).} ``Kaluza-Klein gravity.'' Comprehensive review of KK unification including modern constraints. The KK scalar (dilaton) is the 55-component of the 5D metric and couples conformally to 4D gravity --- structurally identical to $\psi$. \textbf{Grade: A.}

\item \textbf{Alexander et al.\ (2024).} ``Thermodynamic unified field theory: emergent unification.'' Claims unification through thermodynamic principles. Speculative but motivating. \textbf{Grade: C.}
\end{enumerate}

\subsection{DFD's Unification Prospects}

\begin{proposition}[DFD as the Gravitational Sector of a Kaluza-Klein Theory]
\label{prop:KK-embedding}
The DFD field $\psi$ has the same structural role as the Kaluza-Klein dilaton in 5D gravity:
\begin{center}
\begin{tabular}{lcc}
\toprule
& \textbf{KK Dilaton} & \textbf{DFD $\psi$} \\
\midrule
Origin & $g_{55}$ component of 5D metric & Refractive index field \\
Conformal coupling & $e^{2\alpha\phi}g_{\mu\nu}$ & $e^{2\psi}\eta_{\mu\nu}$ \\
EM coupling & $e^{-2\alpha\phi}F_{\mu\nu}F^{\mu\nu}$ & $e^{-2\psi}F_{\mu\nu}F^{\mu\nu}$ (conformal) \\
Kinetic term & $-\frac{1}{2}(\partial\phi)^2$ (canonical) & $\Lambda_\psi^4[\chi-\ln(1+\chi)]$ (MOND) \\
Constraint? & No (propagating) & Yes (constraint field) \\
\bottomrule
\end{tabular}
\end{center}

The key difference is that $\psi$ is a \textbf{constraint field} while the KK dilaton propagates freely. This suggests:
\begin{enumerate}
\item DFD could be a \textbf{constrained} Kaluza-Klein theory, where the extra dimension is ``frozen'' (no KK excitations) and only the zero-mode constraint survives.
\item The MOND kinetic term $\chi - \ln(1+\chi)$ could arise from integrating out the KK tower in a warped extra dimension where the warp factor produces the MOND interpolation function.
\end{enumerate}
\end{proposition}

\begin{theorem}[DFD-Electromagnetism Coupling from Disformal Sector]
\label{thm:EM-coupling}
In DFD, photons couple to $\psi$ through the conformal factor of the physical metric:
\begin{equation}
\mathcal{L}_{\rm EM} = -\frac{1}{4}\sqrt{-\tilde{g}}\,\tilde{g}^{\mu\alpha}\tilde{g}^{\nu\beta}F_{\mu\nu}F_{\alpha\beta}.
\end{equation}

Expanding in the weak-field limit ($\psi \ll 1$):
\begin{equation}
\mathcal{L}_{\rm EM} = -\frac{1}{4}F_{\mu\nu}F^{\mu\nu} - \psi\,F_{\mu\nu}F^{\mu\nu} + D(\chi)\,\partial^\mu\psi\,\partial^\alpha\psi\,F_{\mu\nu}F_\alpha{}^\nu + O(\psi^2).
\end{equation}

The conformal coupling $\psi F^2$ is universal (all massless gauge fields couple this way). The disformal coupling $D\partial\psi\partial\psi F F$ is specific to DFD and \textbf{produces birefringence} of photons in strong $\psi$-gradients:
\begin{equation}
\Delta c_{\rm EM} = D(\chi)(\partial_r\psi)^2 \sim \frac{D_0\chi_r}{(1+\chi_r)^2}\cdot a_\star^2.
\end{equation}

For a galaxy-scale $\psi$-gradient ($\chi \sim 1$): $\Delta c/c \sim D_0 a_0 R_{\rm galaxy}/c^2 \sim 10^{-29}$. Currently undetectable.
\end{theorem}

\subsection{Path to Full Unification}

The honest assessment:
\begin{enumerate}
\item \textbf{DFD unifies gravity + dark matter + dark energy}: confirmed by the single-field structure.
\item \textbf{DFD couples to EM, weak, strong}: through the physical metric $\tilde{g}_{\mu\nu}$. All SM fields are minimally coupled to $\tilde{g}$.
\item \textbf{DFD does not explain the origin of gauge forces}: the SM gauge structure ($SU(3)\times SU(2)\times U(1)$) is input, not derived.
\item \textbf{A KK embedding could potentially unify}: if $\psi$ is the zero-mode of a higher-dimensional metric, the gauge structure could emerge from the isometries of the extra dimensions (Chamseddine-Connes, or standard KK).
\item \textbf{The constraint structure of $\psi$ is a significant obstacle}: KK dilatons propagate freely; $\psi$ is a constraint. A ``frozen KK'' scenario would need a mechanism to impose the constraint from higher-dimensional dynamics.
\end{enumerate}

\subsection{Grade: B}

The KK analogy is structurally exact and well-motivated. The disformal birefringence prediction is new. But full unification remains aspirational --- DFD provides the gravitational sector but does not derive the gauge sector.

\subsection{Hard Question (Agent 15)}

\textbf{If DFD is a constrained KK theory, what mechanism ``freezes'' the extra dimension to produce the constraint structure of $\psi$? In string theory, moduli stabilization freezes extra dimensions, but the moduli become massive propagating fields --- not constraints. Is there a KK stabilization mechanism that produces a constraint field (no propagator) rather than a massive scalar? If so, what does this imply for the string landscape?}

\newpage

%=============================================================================
\section{Summary Table: i34 Experiment Results}
\label{sec:experiment-summary}
%=============================================================================

\begin{center}
\begin{longtable}{lp{6cm}cc}
\toprule
\textbf{Agent} & \textbf{Key Result} & \textbf{Grade} & \textbf{Status} \\
\midrule
\endfirsthead
\midrule
\textbf{Agent} & \textbf{Key Result} & \textbf{Grade} & \textbf{Status} \\
\midrule
\endhead
\bottomrule
\endlastfoot
6 (Measurement) & $\psi$-monitoring provides gravitational einselection; preferred basis = position; MOND-enhanced decoherence & A$-$ & \NEW \\
7 (Entanglement) & $\psi$-vacuum: modified area law with MOND enhancement $1/\mu(\chi)$; no volume law in ground state & B+ & \NEW \\
12 (Experiment) & Two-stage BMV: ground (baseline) + space (MOND, factor 2700 enhancement); definitive DFD/GR discriminator & A & \NEW \\
15 (Unification) & KK-dilaton analogy exact; disformal birefringence predicted; full unification requires ``frozen KK'' mechanism & B & \NEW \\
\end{longtable}
\end{center}

\end{document}
